\subsection{A small split prime detects every extra abelian character}
\label{sec:abelian-character-detector}

The following argument is conditional on Dirichlet GRH. It does not require GRH for a nonabelian extension.
Write $\eta=\chi_{-3}$.

\begin{proposition}\label{prop:abelian-character-detector}
Assume GRH for Dirichlet $L$-functions.
There is an absolute constant $C>0$ with the following property.
Let $\xi$ be a primitive Dirichlet character of conductor $q$,
and suppose $\xi\ne1,\eta$. If
\[
 \xi(p)=1
 \qquad\text{for every rational prime }p\le 2x
       \text{ with }p\equiv1\pmod3,\ p\nmid q,
                                                        \tag{A.1}
\]
where $x\ge2$, then
\[
                         x\le C\log^2(3q).              \tag{A.2}
\]
In particular, $\log(3q)\gg\sqrt{x}$.
The same conclusion holds if the condition in (A.1) is imposed
only for the primes in $[x,2x]$.
\end{proposition}

\paragraph{Proof.}
A nonnegative prime coefficient detects precisely a failure of
$\xi(p)=1$ at a prime split in $\mathbb Q(\sqrt{-3})$.
Its smooth sum has a positive pole term of order $x$.
Dirichlet GRH bounds its zero terms by
$O(\sqrt{x}\log(3q))$, uniformly in the conductor.
Ramified primes and higher powers do not have order $x$.
We supply the uniform zero estimate, including the imprimitive
product $\eta\xi$.

Fix once and for all a nonzero $W\in C_c^\infty((1,2))$ with
$W\ge0$, and write
\[
 \widehat W(s)=\int_0^\infty W(u)u^{s-1}\,du,\qquad
                         w_1=\widehat W(1)>0.
\]
Its Mellin transform decreases faster than every power of $|\Im s|$
on each fixed vertical strip.

We first prove that for a primitive Dirichlet character $\chi$
of conductor $d$, including the character of conductor one,
\[
 \sum_{n\ge1}\Lambda(n)\chi(n)W(n/x)
 ={\bf1}_{\chi=1}w_1x+
             O_W\bigl(\sqrt{x}\log(3d)\bigr).
                                                        \tag{A.3}
\]
For $\chi\ne1$, set $a\in\{0,1\}$ by $\chi(-1)=(-1)^a$ and
\[
 \Lambda_\chi(s)=
 (d/\pi)^{(s+a)/2}\Gamma((s+a)/2)L(s,\chi).
\]
This is entire of order one and satisfies
$\Lambda_\chi(s)=\epsilon_\chi
 \overline{\Lambda_\chi(1-\overline s)}$, with
$|\epsilon_\chi|=1$.
For $\chi=1$ use instead
$\xi_{\mathbb Q}(s)=s(s-1)\pi^{-s/2}\Gamma(s/2)\zeta(s)$.
These analytic properties follow from the Dirichlet case of
Tate's functional equation; see \citet[Theorems~5.16 and 5.22(b),(c), pp.~34--39]{Poonen2015}.

Here is the required local zero count, with its conductor dependence.
The Hadamard product and the functional equation give
\[
 \sum_\rho
 \frac{2-\Re\rho}{(2-\Re\rho)^2+(t-\Im\rho)^2}
 =
 \frac12\log(d/\pi)
 +\frac12\Re\frac{\Gamma'}{\Gamma}\left(\frac{2+a+it}{2}\right)
 +\Re\frac{L'}{L}(2+it,\chi)                              \tag{A.4}
\]
for $\chi\ne1$. Zeros are counted with multiplicity.
To see that no unknown constant occurs, differentiate the
genus-one Hadamard product. Its real part is the displayed zero
sum plus a real constant. On $\Re s=1/2$ the functional equation
makes the logarithmic derivative purely imaginary. The zero
terms there cancel in the pairs $\rho,1-\overline\rho$;
under the granted GRH each individual real part is zero.
The constant is therefore zero. The sums of real parts converge
absolutely, since the function has order one and $0<\Re\rho<1$.
For $\xi_{\mathbb Q}$ the identity has the additional terms
$\Re(1/(2+it)+1/(1+it))$, which are bounded.

At $\Re s=2$, the logarithmic derivative of the Euler product
is bounded absolutely by $\sum_n\Lambda(n)n^{-2}$.
The gamma logarithmic derivative is
$O(\log(3+|t|))$ uniformly on the indicated line.
Every zero with $|\Im\rho-t|\le1$ contributes at least $1/5$
to the left side of (A.4). Consequently
\[
 \#\{\rho:|\Im\rho-t|\le1\}
       \ll \log\bigl(3d(3+|t|)\bigr).                    \tag{A.5}
\]
All constants here are absolute.
This also proves
\[
 \sum_\rho|\widehat W(1/2+i\Im\rho)|
                        \ll_W \log(3d),                 \tag{A.6}
\]
by summing (A.5) over unit height intervals and using Mellin
decay of order, for example, four.

Mellin inversion of $-L'/L$ on $\Re s=2$ and contour displacement
to $\Re s=-1/2$ give
\[
 \sum_{n\ge1}\Lambda(n)\chi(n)W(n/x)
 ={\bf1}_{\chi=1}w_1x
   -\sum_\rho x^\rho\widehat W(\rho)+O_W(\log(3d)).
                                                        \tag{A.7}
\]
For completeness, the displaced integral is
$O_W(x^{-1/2}\log(3d))$: the functional equation expresses
$L'/L(-1/2+it,\chi)$ using the absolutely convergent Euler
product at $3/2-it$ and the gamma logarithmic derivatives.
The possible trivial zero at zero contributes $O_W(1)$.
Choose horizontal heights in successive unit intervals a distance
$\gg1/\log(3d(3+|t|))$ from all zero ordinates, using (A.5).
The logarithmic derivative obtained from the Hadamard product
has polynomial logarithmic growth on those contours, and the
rapid Mellin decay makes their integrals tend to zero.
This justifies the displacement and absolute convergence of its
zero sum. Under GRH, $|x^\rho|=\sqrt{x}$.
Equations (A.6) and (A.7) imply (A.3).

If a character modulo $m$ is induced from a primitive character
of conductor $d\mid m$, replacing it by that primitive character
changes its smooth prime-power sum only at $p\mid m$.
For a fixed $p$, the number of exponents $k$ with
$x<p^k<2x$ is at most $1+\log2/\log p$. Hence the total change is
\[
 \ll_W \sum_{p\mid m}\log p
 \ll_W\log(3m).                                         \tag{A.8}
\]
Thus (A.3), with $d$ replaced by $m$, also holds for an
imprimitive character, retaining its correct principal pole term.

Define, on prime powers,
\[
 h(p^k)=\frac{1+\eta(p)^k-\Re(\xi(p)^k)
                    -\Re((\eta\xi)(p)^k)}2,
\]
where the usual zero extension is used at primes dividing
a character's modulus. Equivalently use the four Dirichlet
characters $1,\eta,\xi,\eta\xi$ at the integer $p^k$.
The last three are nonprincipal: the assumptions
$\xi\ne1,\eta$ exclude both possible principal products.
The primitive conductor inducing $\eta\xi$ divides $3q$.
Applying (A.3) and (A.8) therefore gives
\[
 \sum_{p,k\ge1}(\log p)h(p^k)W(p^k/x)
       =\frac{w_1x}{2}
                       +O_W(\sqrt{x}\log(3q)).           \tag{A.9}
\]
For $p\nmid3q$,
\[
 h(p)=\frac{(1+\eta(p))(1-\Re\xi(p))}{2}\ge0.
\]
It is zero at inert primes and, by (A.1), at every split
prime in the support of $W(p/x)$. The exceptional prime terms
with $p\mid3q$ contribute $O_W(\log(3q))$.

Finally $|h(p^k)|\le2$, and the elementary Chebyshev bound
$\vartheta(y)=\sum_{p\le y}\log p\ll y$ yields
\[
 \sum_{\substack{k\ge2\\p^k\le2x}}\log p
 \le\sum_{2\le k\le\log_2(2x)}\vartheta((2x)^{1/k})
 \ll \sqrt{x}+x^{1/3}\log(3x)
 \ll\sqrt{x}.                                           \tag{A.10}
\]
The last inequality holds with an absolute constant for
all $x\ge2$. Inserting these bounds into (A.9) gives
$w_1x/2\ll_W\sqrt{x}\log(3q)$, proving (A.2).

\paragraph{Conductor--discriminant consequence.}
If $F/\mathbb Q$ is abelian and contains
$K=\mathbb Q(\sqrt{-3})$, and every degree-one prime of $K$
above a rational prime $p\le2x$ splits completely in $F/K$,
then every character in the character group of $F/\mathbb Q$
other than $1,\eta$ has conductor $q_\xi$ satisfying
$\log(3q_\xi)\gg\sqrt{x}$.
By the conductor--discriminant formula,
\[
 \log\operatorname{rd}(F)
   =\frac1{[F:\mathbb Q]}\sum_\xi\log q_\xi.
\]
For $F\ne K$, its degree is an even integer at least four,
so at least half its characters are outside $\{1,\eta\}$.
It follows, for all sufficiently large $x$, that
$\log\operatorname{rd}(F)\gg\sqrt{x}$, with an absolute
constant. This last conclusion needs the stated complete
splitting condition; preserving a selected trace through
some other mechanism is not asserted to imply that condition.

